\section{Scope, evidence, and how to read the review}
The central object is a model that predicts how a controlled robot evolves. The
proposed transfer problem changes the contact geometry after pretraining that
model on a motion library. A successful research result would explain both the
additional data needed and the reliability of the policy updates obtained from
the adapted model.

The eleven papers were selected because each isolates a necessary part of that
question. PETS concerns uncertainty and planning; MBPO concerns how much model
rollout to trust; Dreamer concerns learned latent dynamics and imagination;
ContactNets concerns contact structure; SHAC concerns short simulator-gradient
windows. PWM, DMO, RWM-U, SuperMimic, $\nabla$Mimic, and WM-LOCO make the more direct
connections to learned gradients, humanoid motion, and terrain-aware control.

This is a question-driven selection, not an exhaustive survey. General video
generation, language-conditioned robot foundation models, and broad motion
retargeting surveys are outside its scope. A model that predicts photorealistic
video is not automatically a model with adequate action sensitivity for motor control.

Each paper is treated using the same sequence: its role, a concise account of
the reported method, a normalized mathematical skeleton, an analysis of what
that skeleton means, and a test or limitation relevant to contact transfer.
The analysis is original to this review; it should not be mistaken for an
experiment performed by the cited authors.

\begin{limitation}{Evidence discipline}
No reviewed result establishes the number of interactions required for our exact
flat-motion-to-new-stair-geometry experiment. That quantity depends on the
source data, target distribution, observations, controller, model class, and
success criterion. A reported wall-clock speedup on motion tracking cannot
substitute for this missing measurement.
\end{limitation}

Read Sections~\ref{sec:notation}--\ref{sec:contactproblem} first. For the most
direct path, continue with PWM, DMO, SuperMimic, $\nabla$Mimic, and WM-LOCO.
The earlier papers explain the model-usage and representation choices behind
them. Sections~\ref{sec:comparison}--\ref{sec:experiments} turn the review into
an experimental question.

\section{One notation for all methods}\label{sec:notation}
\subsection{Physical state, environment, task, and observations}
Let $x_t$ be the physical state of the robot and any moving objects included in
the system. Static environmental geometry is $E$; material, calibration, and
other physical parameters are $\xi$. The transition is
\begin{equation}\label{eq:true}
x_{t+1}\sim P_{E,\xi}(\cdot\mid x_t,a_t),\qquad
x_{t+1}=\Fsim(x_t,a_t;E,\xi)\quad\text{when deterministic}.
\end{equation}
Here $a_t$ is the action actually executed, after any exploration noise and
clipping. An explicit actuation map $u_t=\mathcal U(a_t)$ converts it to a physical
setpoint or torque when needed. Numerical comparisons must specify whether
action coordinates are normalized values, radians, or torques.
If an implementation predicts an offset relative to a reference, we include
the reference addition inside the policy and call its resulting executed command
$a_t$. The model must receive the command that actually determines actuation;
an offset alone is insufficient when its reference-dependent baseline is unknown.
$\mathcal U$ is an actuator map, while $M(q)$ below always denotes a mass matrix.

Write $\Gtask_t$ for the complete task/reference specification, including its
phase where relevant. The supplied policy-goal vector is
$g_t=G(x_t,\Gtask_t)$. This distinction matters: a compact actor goal can omit
reference quantities still used by the training reward. We associate reward
$r_t$ with the transition $t\rightarrow t+1$:
\begin{equation}
r_t=R(x_t,a_t,x_{t+1};\Gtask_t,E),\qquad
o_t\sim O(\cdot\mid x_t,E,\xi).
\end{equation}
Papers that call this reward $r_{t+1}$ have simply shifted the label in this review.
The observation/action history is $\Hhist_t$. A geometry query or perceptual
encoder supplies local environment information $e_t$:
\begin{equation}
e_t=\mathcal Q(E,x_t)\ \text{with known geometry},\qquad
e_t=\mathcal E_\phi(\Hhist_t)\ \text{when inferred}.
\end{equation}
The same terrain map must be queried at the predicted pose during imagination.
Using future simulator observations inside a learned rollout would change the
information available to the algorithm.

\subsection{Network roles and probability conventions}
\begin{longtable}{@{}L{0.22\linewidth}L{0.73\linewidth}@{}}
\toprule Symbol & Meaning throughout this review\\\midrule\endhead
$\theta$ & Policy parameters, including task encoders/projections assigned to the policy.\\
$\phi$ & Learned dynamics, observation encoders, and predictive heads.\\
$\psi$ & Value-function parameters; $\bar\psi$ denotes a target copy.\\
$\policy$ & Policy, deterministic or reparameterized stochastic as specified.\\
$\Fhat$ & Learned physical-state transition or transition mean.\\
$\widehat P_\phi$ & Learned transition distribution.\\
$z_t=(m_t,\ell_t)$ & Learned latent state: deterministic memory and stochastic component when present.\\
$\mathsf q_\phi$ & Observation-conditioned posterior over a latent state.\\
$\mathsf p_\phi$ & Latent transition prior, propagated without the next observation.\\
$\widehat R_\phi$ & Learned reward predictor, when the algorithm uses one.\\
$\zeta_t$ & Learned motion-goal embedding, distinct from actuator command $u_t$.\\
$V_\psi,Q_\psi$ & State-value and action-value functions.\\
$\gamma,\lambda$ & Discount and lambda-return mixing coefficient.\\
$c_t,C_t$ & One-step continuation and cumulative continuation, $C_t=\prod_{i<t}c_i$.\\
$\sg(\cdot)$ & Stop-gradient: preserve a value but suppress that derivative path.\\
$H,K$ & Policy-optimization horizon and model-fitting horizon, respectively.\\
$N_{\rm pre},N_{\rm new}$ & Pretraining and additional target-environment transitions.\\
\bottomrule
\end{longtable}
Parameter labels denote roles, not a requirement that implementations have
three disjoint optimizers. An actor and critic may share a trunk. The policy's
input may be $x_t$, a history, or $z_t$; using the same $\theta$ does not erase
those architectural differences.

We always call an observation-conditioned latent distribution a posterior
$\mathsf q_\phi$ and a prediction-only distribution a prior $\mathsf p_\phi$.
This reverses the letters used by the original Dreamer paper, intentionally;
the distinction is stated again in its method section. Contact impulses are
$j$, not $\lambda$, so they cannot be confused with the return coefficient.

\subsection{State models versus latent models}
A physical-state model uses
\begin{equation}
\hat x_{t+1}=\Fhat(\hat x_t,a_t,e_t,\xi).
\end{equation}
A latent model instead infers a posterior from observations and rolls a prior:
\begin{align}
z_t&\sim\mathsf q_\phi(z_t\mid z_{t-1},a_{t-1},o_t),\\
\hat z_{t+1}&\sim\mathsf p_\phi(z_{t+1}\mid\hat z_t,a_t).
\end{align}
Some inputs are absent in a particular paper. We never infer that a method uses
terrain conditioning merely because $e_t$ appears in this general template.
A task embedding, an environment encoding, and a belief over unobserved
physical parameters are three different kinds of conditioning.

\section{A common mathematical lens for learning and control}\label{sec:common}
\subsection{Fitting dynamics}
A physical model can be trained by negative log likelihood,
\begin{equation}\label{eq:nll}
L_{\rm dyn}(\phi)=-\E_{\mathcal D}\log\widehat P_\phi(x_{t+1}\mid x_t,a_t,e_t,\xi),
\end{equation}
or a multistep residual loss. With $\hat x_0=x_0$ and logged actions,
\begin{equation}\label{eq:multistep}
L_{\rm dyn}^{(K)}=\E_{\mathcal D}\sum_{k=0}^{K-1}w_k
\norm{\Delta(x_{k+1},\hat x_{k+1})}_{W}^{2},\quad
\hat x_{k+1}=\Fhat(\hat x_k,a_k,e_k,\xi).
\end{equation}
$\Delta$ must respect the state representation: quaternion coordinates should
not be treated as unconstrained Euclidean angles. A latent model may instead
learn observation/reward likelihoods and a prior--posterior consistency term.

These objectives measure predictions on a data distribution. They do not
directly certify errors at policy-chosen states outside that distribution or
derivatives of those predictions.

\subsection{Short-horizon actor--critic as the common backbone}
Let $y_t$ temporarily mean the optimization state: $x_t$ for a physical-state
model and $z_t$ for a latent model. This is a placeholder, not an extra latent.
Many methods use a variant of
\begin{equation}\label{eq:shortobjective}
J_H=\E\left[\sum_{t=0}^{H-1}\gamma^t C_t\widetilde r_t
+\gamma^H C_HV_\psi(y_H,g_H)\right],
\end{equation}
where $\widetilde r_t$ can be an analytic reward or a learned prediction.
Continuation can be hard termination or a soft surrogate. Reference-end masks
are an additional implementation detail, not implied by this notation.

One useful normalized value target is
\begin{align}\label{eq:lambda}
Y_H&=V_\psi(y_H,g_H),\\
Y_t&=\widetilde r_t+\gamma c_t\big[(1-\lambda)V_\psi(y_{t+1},g_{t+1})+\lambda Y_{t+1}\big],\\
L_V&=\E\big[V_\psi(y_t,g_t)-\sg(Y_t)\big]^2.
\end{align}
This template does not claim identical target networks, reward normalization,
discount indexing, or terminal handling in every method. Such differences are
identified in the paper analyses.

\subsection{Four different ways a model can enter control}
\begin{enumerate}
\item \textbf{Planning}: score candidate actions under model rollouts and select
one, without differentiating a neural policy through the model.
\item \textbf{Imagined policy optimization}: generate predicted trajectories and
differentiate their return with respect to policy parameters.
\item \textbf{Decoupled forward/backward}: use physical-simulator states in the
forward pass and learned local Jacobians in the backward pass.
\item \textbf{Predictive features}: feed a learned recurrent representation to
a policy trained by another estimator, such as PPO.
\end{enumerate}
These choices should be compared separately from the choice of state versus
latent representation. ``Uses a world model'' is not a complete algorithm description.

\begin{figure}[htbp]
\centering
\begin{tikzpicture}
\node[box,text width=33mm] (a) at (0,0) {Policy\\action};
\node[box,text width=33mm] (b) at (4.8,0) {Learned transition\\predicted state};
\node[box,text width=33mm] (c) at (9.6,0) {Predicted return\\plus value tail};
\draw[arr](a)--(b);\draw[arr](b)--(c);
\draw[grad](c.south)--++(0,-0.6)-|node[pos=0.25,below,font=\scriptsize]{Model derivatives through imagined states}(a.south);
\node[box,text width=33mm] (d) at (0,-2.4) {Policy\\action};
\node[box,text width=33mm] (e) at (4.8,-2.4) {Simulator\\actual state};
\node[box,text width=33mm] (f) at (9.6,-2.4) {Simulator return\\plus value tail};
\draw[arr](d)--(e);\draw[arr](e)--(f);
\draw[grad](f.south)--++(0,-0.6)-|node[pos=0.25,below,font=\scriptsize]{Learned derivatives evaluated at actual states}(d.south);
\node[box,text width=33mm] (g) at (0,-4.8) {Policy\\action};
\node[box,text width=33mm] (h) at (4.8,-4.8) {Differentiable\\simulator};
\node[box,text width=33mm] (i) at (9.6,-4.8) {Simulator return\\plus value tail};
\draw[arr](g)--(h);\draw[arr](h)--(i);
\draw[grad](i.south)--++(0,-0.6)-|node[pos=0.25,below,font=\scriptsize]{Simulator adjoint with its stated derivative rules}(g.south);
\node[box,text width=33mm] (j) at (0,-7.2) {Learned predictive\\memory};
\node[box,text width=33mm] (k) at (4.8,-7.2) {Policy\\action};
\node[box,text width=33mm] (l) at (9.6,-7.2) {Simulator rewards\\and advantages};
\draw[arr](j)--(k);\draw[arr](k)--(l);
\draw[grad](l.south)--++(0,-0.6)-|node[pos=0.25,below,font=\scriptsize]{PPO score-function update}(k.south);
\end{tikzpicture}
\caption{Four policy-learning routes. Solid arrows show forward information;
dashed arrows label the update estimator. The final row uses advantages and
action log probabilities, not differentiation of simulator reward through the
dynamics. Inputs and repeated rollout steps are omitted for readability.}
\end{figure}

\subsection{A derivative is an additional modeling requirement}
For a deterministic policy at a fixed task schedule, suppress other inputs and
write $a_t=A_\theta(x_t)$. The pathwise sensitivity is
\begin{equation}\label{eq:sensitivity}
D_\theta x_{t+1}=(F_x+F_aA_x)D_\theta x_t+F_aA_\theta.
\end{equation}
The matrices $F_x,F_a$ come from a simulator or a learned function. History,
latent memory, and state-dependent goal construction can be included by
augmenting the optimization state. Freezing $\phi$ during the actor update
does not remove $\partial\Fhat/\partial a$ or $\partial\Fhat/\partial x$.
When $e=\mathcal Q(E,x)$, the total state derivative also includes the geometry
query path where that path is differentiable. A discrete query requires a
declared piecewise derivative or surrogate; silently dropping it changes the update.

This equation identifies two independent failure paths: the forward states can
be wrong, or the local sensitivities can be wrong even at good forward states.
Products of sensitivities can amplify errors as $H$ grows.

\begin{interpretation}{Original illustration: small prediction error, large derivative error}
Consider $F(a)=a$ and $\widehat F(a)=a+\epsilon\sin(\omega a)$. Prediction error
is at most $\epsilon$, yet derivative error can reach $\epsilon\omega$. Taking
$\epsilon=0.01$ and $\omega=200$ gives at most 0.01 state error and derivative
error as large as two. The predicted derivative can even point in the opposite
direction. This elementary construction is not an empirical result from a
reviewed paper; it explains why a trajectory-fit metric alone cannot certify
useful policy gradients.
\end{interpretation}
\begin{figure}[htbp]\centering
\includegraphics[width=\linewidth]{figures/prediction_and_gradient.pdf}
\caption{Synthetic review illustration of the distinction between prediction
and derivative errors. No paper measurements are plotted.}
\end{figure}

\section{What changes when the contact geometry changes?}\label{sec:contactproblem}
\subsection{The missing-input problem}
For two environments $E_1,E_2$, identical $(x,a)$ can yield
\begin{equation}
\Fsim(x,a;E_1,\xi)\ne\Fsim(x,a;E_2,\xi).
\end{equation}
A model with no environment information cannot distinguish those outcomes in
general. A larger network cannot recover information that its inputs omit.
With one fixed known map and suitable global coordinates, location can encode
geometry implicitly; that is not generalization to arbitrary unseen maps.

Flat-ground motion already involves impacts, support changes, slipping, and
body contacts. The proposed transfer changes \emph{which surfaces exist and
where they are}, potentially also which body parts touch them. Later tasks may
change friction, compliance, object motion, or contact laws as well. Those are
distinct axes and should not be changed simultaneously in the first experiment.

\subsection{An idealized contact template}
For configuration $q$ and generalized velocity $v$, a rigid-contact template is
\begin{align}\label{eq:mechanics}
M(q)(v_{t+1}-v_t)&=\Delta t\,f_{\rm free}(x_t,u_t)+J(q,E)^Tj_t,\\
q_{t+1}&\approx q_t+\Delta t\,\Gamma(q_t)v_{t+1},\\
d_i(q_{t+1},E)&\ge0,\quad j_i^n\ge0,\quad d_i(q_{t+1},E)j_i^n=0,\\
\norm{j_i^\tau}&\le\mu_i j_i^n.
\end{align}
$d_i$ is signed distance, $J$ maps generalized to contact-frame velocity, and
$j$ is impulse. $\Gamma$ handles configuration representations such as floating-base
quaternions. Complementarity says contact pushes only at an active surface.
This is an explanatory rigid-contact template, not a claim that every reviewed
simulator uses exactly this forward law; compliant contact models differ.

When contact identities or modes change, a derivative taken with the current
active set may not explain the effect of a finite action change that creates a
different contact. Simulator exactness, hardware fidelity, and usefulness for
optimization are consequently separate questions.

\subsection{What a terrain-aware imagined rollout must maintain}
If $e_t$ is an egocentric map, advancing $\hat x_t$ requires advancing the map
query or predicting the relevant latent scene. Holding an arbitrary local
encoding constant while the robot moves can give inconsistent predictions.
For moving objects, their states must be modeled too. For hidden materials or
occluded geometry, a history-dependent belief may be necessary.

The task must also remain coherent. Tracking a world-aligned flat-ground
reference while climbing may penalize the correct height increase. Terrain-aware
dynamics alone does not repair a mismatched goal, reward, critic input, or
termination rule. These distinctions guide the paper selection that follows.
